\section{A noisy Markov chain is a tree}
\label{sec:tree}

\begin{figure}[t]
  \centering
  \begin{tikzpicture}[
    var/.style={circle,draw,thick,minimum size=8.2mm,inner sep=0pt,font=\small},
    obs/.style={circle,draw,thick,fill=black!12,minimum size=8.2mm,inner sep=0pt,font=\small},
    fac/.style={rectangle,fill=black,minimum size=2.2mm,inner sep=0pt},
    lab/.style={font=\small,inner sep=1pt},
    msg/.style={-{Stealth[length=2mm]},thick}]
  \node[var] (a1) at (0,0) {$a_1$};
  \node (d1) at (1.15,0) {$\cdots$};
  \node[var] (ak) at (2.3,0) {$a_k$};
  \node[var] (ak1) at (4.0,0) {$a_{k+1}$};
  \node (d2) at (5.15,0) {$\cdots$};
  \node[var] (aL) at (6.3,0) {$a_L$};
  \node[fac] (m0) at (-0.8,0) {};
  \draw (m0) -- (a1);
  \node[lab,above=0.8mm of m0] {$\mu$};
  \foreach \u/\v/\n in {a1/d1/ka,d1/ak/kb,ak/ak1/kc,ak1/d2/kd,d2/aL/ke} {
    \node[fac] (\n) at ($(\u)!0.5!(\v)$) {};
    \draw (\u) -- (\n) -- (\v);
  }
  \node[lab,above=0.8mm of kc] {$K$};
  \foreach \i/\n in {a1/1,ak/k,ak1/{k+1},aL/L} {
    \node[fac] (o\i) at ($(\i)+(0,-0.9)$) {};
    \node[obs] (x\i) at ($(\i)+(0,-1.8)$) {$x_{\n}$};
    \draw (\i) -- (o\i) -- (x\i);
  }
  \node[lab,left=0.8mm of oak] {$\ell_t$};
  \draw[msg] ($(a1.north)+(-0.25,0.5)$) -- node[lab,above]{$f_k$} ($(ak.north)+(-0.1,0.5)$);
  \draw[msg] ($(aL.north)+(0.25,0.5)$) -- node[lab,above]{$b_k$} ($(ak.north)+(0.1,0.5)$);
\end{tikzpicture}

  \caption{Factor graph of the posterior $p_t(\mathbf a\mid\mathbf x)$ of \cref{eq:posterior_chain}. White circles are
  the clean values $a_k$, grey circles the noisy observations $x_k$, which are known. Black squares are the factors: the
  law $\mu$ of the first site, the kernel $K(a_{k+1}\mid a_k)$ between neighbours and the noise term
  $\ell_t(x_k\mid a_k)$ at each site. There is no loop, so belief propagation is exact. The forward message $f_k$ carries
  what $x_1,\ldots,x_k$ say about $a_k$, the backward message $b_k$ what $x_{k+1},\ldots,x_L$ say about it, and the
  posterior of $a_k$ is proportional to $f_k(a_k)\,b_k(a_k)$.}
  \label{fig:factor-graph}
\end{figure}

\paragraph{Beyond Gaussianity.}
The Gaussian chain is exactly solvable and linear. We now take a controlled non-Gaussian step, the Laplace chain
\begin{equation}
    a_{k+1}=\alpha a_k+\varepsilon_k,
    \qquad
    \varepsilon_k\sim\operatorname{Laplace}(0,b),
    \label{eq:laplace_chain}
\end{equation}
with $2b^2=1-\alpha^2$. The noise has the same variance as in the Gaussian chain, so the two chains have the same
covariance. But the Laplace noise has a sharp peak and heavier tails, and the score is no longer linear in $\mathbf x$.

\paragraph{The posterior factor graph.}
The first-order Markov structure survives the noise. Combining the kernel with the independent noise terms gives
\cref{eq:posterior_chain}, whose factor graph is a chain with one extra factor on each site (\cref{fig:factor-graph}).
Since the noise acts on each site separately, conditioning on the observations adds no edge. The posterior is still
connected and has no loop: it is a tree, where the cavity method and belief propagation are exact
\cite{mezard2001bethe,mezard2006reconstruction,mezard2009information}.

\paragraph{Exact inference by belief propagation.}
On a tree, the sum--product algorithm gives the exact marginals with one forward and one backward pass
\cite{pearl1988probabilistic,kschischang2001factor,yedidia2003understanding}. With forward messages $f_k$ and backward
messages $b_k$, the recursions are
\begin{align}
    f_1(a_1)
    &=\mu(a_1)\ell_t(x_1\mid a_1),\nonumber\\
    f_{k+1}(a_{k+1})
    &=\ell_t(x_{k+1}\mid a_{k+1})\nonumber\\
    &\qquad\times\int K(a_{k+1}\mid a_k)f_k(a_k)\,da_k,\nonumber\\
    b_L(a_L)&=1,\nonumber\\
    b_k(a_k)
    &=\int K(a_{k+1}\mid a_k)
      \ell_t(x_{k+1}\mid a_{k+1})\nonumber\\
    &\qquad\times b_{k+1}(a_{k+1})\,da_{k+1}.
    \label{eq:bp}
\end{align}
The posterior of one site is then
\begin{equation}
    p(a_k\mid\mathbf x)
    \propto f_k(a_k)b_k(a_k),
\end{equation}
and its mean gives the score through \cref{eq:tweedie}:
\begin{equation}
    S_{t,k}(\mathbf x)=\frac{e^{-t}}{\Delta_t}\,
    \frac{\int a\,f_k(a)b_k(a)\,da}{\int f_k(a)b_k(a)\,da}-\frac{x_k}{\Delta_t}.
\end{equation}

So, for any known first-order kernel, two sweeps of messages give the exact joint score, with no assumption of
Gaussianity or linearity. The noise level only enters through $\ell_t$, so one kernel gives the score at every $t$.
The same exact inference on a tree is at the heart of the hierarchical models of \cite{garnierbrun2024transformers},
where transformers trained on such data learn to approximate BP; a Markov chain is the simplest tree.
